% =============== Related Work =================================================
% Goal: 1 page · structure:
%   (1) Memory plugins (mem0/Letta/Zep) - retrieval focus
%   (2) White-box persona monitoring (Persona Vectors)
%   (3) Cross-encoders & rerankers
%   (4) LLM agent benchmarks
%   (5) Anchor-based behavior detection (analogies)
% =============================================================================

\subsection{Memory Systems for LLM Agents}
\label{sec:related-memory}

A growing line of work addresses the limited context window of LLMs
through external memory. \textbf{mem0} \citep{mem0} stores
key-value memories extracted from conversations and retrieves them
via dense embeddings; the system targets production agent use cases
and reports retrieval Recall@5 in the 0.5--0.6 range on
LongMemEval-S. \textbf{Letta} (formerly MemGPT) \citep{letta}
introduces a hierarchical memory architecture with archival
storage and core memory, modeling LLM context management as an
operating-system problem. \textbf{Zep} \citep{zep} adds temporal
knowledge-graph structure on top of dense retrieval, optimizing for
multi-session temporal reasoning. \textbf{A-MEM} \citep{amem}
introduces dynamic links between memory entries with supersede
detection, addressing memory staleness.

These systems share a common assumption: \emph{the LLM will honor
retrieved memory once it is injected into the context}. We argue
that this assumption fails systematically over long sessions, and
that complementary mechanisms are needed to monitor whether the
agent's behavior actually aligns with the retrieved guidance.

\subsection{White-box Persona Monitoring}
\label{sec:related-persona}

\citet{personavectors2025} represents the state of the art in
white-box persona monitoring: they identify linear directions in
the activation space of an LLM corresponding to traits such as
sycophancy, hallucination propensity, and harmful intent. By
projecting hidden activations onto these directions during
inference, the authors demonstrate the ability to monitor and
steer trait expression at deployment time, predict trait shifts
induced by finetuning, and flag training data likely to cause
undesirable persona changes.

The chief limitation, from the perspective of the typical
end-user of a coding agent, is that this method requires access
to model weights or hidden states. This puts it out of reach of
users interacting with proprietary LLMs (Claude, GPT-4, Gemini)
through API or product surfaces. Our work targets exactly this
deployment gap: providing a black-box analog that runs in
user-space hooks without model-internal access.

\subsection{Embedders and Cross-Encoder Rerankers}
\label{sec:related-embed}

We build on the BGE family of embedding models \citep{bgem3}.
BGE-m3 in particular offers multilingual support across 100+
languages with a 1024-dimensional output, making it suitable for
mixed Chinese/English production agent contexts. For retrieval
reranking, we leverage \texttt{bge-reranker-v2-m3}
\citep{bgereranker}, a cross-encoder trained on similar data,
which scores (query, candidate) pairs jointly rather than
embedding them independently. Foundational sentence embedding
methodology follows \citet{sbert}.

A natural question is whether a cross-encoder can substitute for
the bi-encoder cosine in our drift detection step (anchor
matching). In Section~\ref{sec:eval-rerank-null}, we report a
\emph{negative} result: cross-encoder substitution improves drift
AUC by only 0.0008 at 36$\times$ the latency, suggesting that the
short, semantically heterogeneous nature of behavioral anchor
text leaves little token-interaction signal for cross-encoders to
exploit, in contrast to the long-document setting of retrieval
reranking.

\subsection{Long-Context Memory Benchmarks}
\label{sec:related-bench}

We evaluate retrieval performance using LongMemEval-S
\citep{longmemeval2024}, which provides 500 question-haystack
pairs across six question types: knowledge update, multi-session
synthesis, single-session-user/assistant/preference factual
recall, and temporal reasoning. Each haystack contains roughly 50
candidate sessions, making top-5 retrieval a meaningful target
metric. Other relevant benchmarks include LoCoMo (long
conversation memory) and PerLTQA (Chinese long-term persona
dialogue), which we leave for future evaluation.

\subsection{Anchor-Based Behavior Detection (Analogies)}
\label{sec:related-anchors}

Our anchor-based approach draws conceptual lineage from several
threads. \textbf{Reward modeling in RLHF} aggregates many
preference comparisons into a scalar reward; in our case we
aggregate cosine similarity to many positive/negative behavioral
anchors into a scalar drift score. \textbf{Behavioral cloning}
in robotics learns from demonstrations; our positive anchors
function as ``demonstrations of correct task patterns''.
\textbf{Few-shot in-context learning} relies on a small set of
examples to bias the LLM's distribution; our negative anchors
function as ``examples of behaviors to avoid'', albeit operating
externally rather than in-context.

The core innovation in our system is not any single technical
component but the empirical methodology for designing anchor sets
that yield strong detection signal (Section~\ref{sec:method-evolution}),
together with the integration with hook-level production
deployment.

\subsection{Strategy Distillation}
\label{sec:related-strategy}

\citet{dptagent} introduces distillation path-triggered (DPT)
agents, which extract reasoning patterns from past task
completions and inject them when similar tasks recur. We
incorporate a lightweight DPT-style mechanism for cross-session
strategy retrieval (Section~\ref{sec:method-strategy}), though it
is not the focus of this paper. Our drift detection is
orthogonal to and compatible with strategy distillation.
